
% 31 Mar: Clean up

\subsection{Installment 3: “Folkebladet” March 27, 1918}


% Sverdrup on the Congregation’s Proper Form


The Free Church says only this: the congregation is divinely instituted. No, say the others.

It is precisely here that the Free Church’s kinship with the Congregationalists shows itself. What Congregationalists call their ‘distinctive function’ is “to reveal and to realise the true idea of the church.” In the work of an English writer we read the following:

\begin{quote}
Dr. R. W. Dale, in an essay written ..... years ago, closes with a statement of what he considered to be ‘the distinctive function’ of Congregationalists.

It was he thought, not the work of Evangelisation, for that is the work alike of all Christians; still less was it the work of testifying against the evils, real or supposed, of other churches—for example, the evil of ‘ecclesiastical establishment’; nor was it the work of reconstructing theological science, or taking a chief place in the controversy with unbelief. All such work may form a part of their duty, and sometimes an inevitable part. But their distinctive function is ‘to reveal and to realise the true idea of the Church.’
\end{quote}

One must not misunderstand such an analogy (or one like it). Such a misunderstanding seems to be contained in the “Kristelig Folkekalender” for 1917, page 43, where it is said: “The mission of the Free Church, then, is not first and foremost to uphold the Congregationalist principle of church order among us Norwegian Lutherans, but rather to work for the renewal and liberation of the congregations.”

This is no doubt meant as a correction addressed to me. But I have not ascribed that task to the Free Church. I have only said that the Free Church’s distinctive feature among Lutherans lies in the area of church order, and that its principles of polity resemble those of the Congregationalists very closely, which Professor J. O. Andersen at the University of Copenhagen also maintains in his article on Sverdrup in “Kirkeleksikon for Norden” (three columns in issue no. 66, year 1917): “The new C. (the Free Church) was ordered according to decidedly Congregationalist principles.” I should add that Andersen here is speaking of the Principles and Rules, and that he has the ‘Collected Writings’ before him. Further, that Andersen, as author of a larger work on church organization, understands tolerably well what Congregationalist principles are.

Distinctiveness, I said; not task. There is a difference between these two things. There are, for example, one hundred students at a school. One of them is black; the others are white. There is something distinctive about the one: he is black. But we will not say that it is first and foremost his task to be black.

The distinctive feature of the Free Church among Lutherans lies in the area of church order. But since it has essentially the same principles of polity as the Congregationalists, it may also be said of it that it has the same “distinctive function.” For the Free Church is not a congregation, but an association of congregations which takes on tasks that the individual congregation cannot well carry out on its own.

The distinctive function of the Free Church is to bring the congregation forth in its proper form, or “to reveal and realise the true idea of the church.” Therefore it declares as its purpose “to labor for the enlivening and liberation of all Norwegian Lutheran congregations, so that according to their calling and ability they may work in spiritual independence and freedom for the cause of God’s kingdom at home and abroad .....” (Rule 2).

And how is this purpose to be achieved? Is it by constantly preaching awakening? Rule 3 says: “This purpose is to be achieved chiefly by the training of men and women for spiritual labor in and for the congregation, by the holding of greater and smaller meetings, by the dissemination of suitable literature, by the organization of Free-Church committees and associations, by sending out workers, and by other means which from time to time may show themselves expedient.”

Nota bene—workers as well. Pastor Nisløv is no doubt not entirely in agreement with Rule 3. For he says in the previous issue of this paper: the work of emissaries “will become a foreign plant in the Free-Church soil.”

Now the editors will scarcely be among those who compare the distinctive function of the Free Church with that of the Congregationalists. And yet it seems that what the Congregationalists call their distinctive function, “to reveal and to realise the true idea of the church,” becomes a main matter, indeed the main matter, for the editors. “The main matter for us is that the congregation come forth in its proper form.” (Folkebladet no. 34.)

And yet we read this: “We know, however, that according to God’s Word there is a congregation where two or three are gathered in Jesus’ name .... or where two or more believers are gathered around God’s Word and the sacraments. The fact that a group of people forms an association does not create any congregation, but the fact that souls come to life in God. For us in the Free Church, who have this view, the organization is of less importance.” (The editors.)

Here it is suggested that the main thing is that souls come to life in God, whereas just before it is said that the main thing is that the congregation come into its proper form. Presumably this is meant to signify the same thing. If it does not do so for the editors, then it does so for many in the Free Church.

But this is not the view of the Principles. In his explanation of Principle No. 1, S. says: “And next to the question of the salvation of the individual soul, there is no more important question than that of the congregation.”

Thus, two questions.

Now we know that the first is included in the second; but the second is not simply included in the first. We know also that Sverdrup used the word awakening in a very comprehensive sense: the gifts of grace too were to be awakened. And therefore he could say: “That which is really the main thing in the free congregation, and which we cannot stress often enough, is the unfolding and right use of the gifts of grace” (C. W. II).

That souls come to life in God does not without more ado mean that the gifts of grace are unfolded and put to use, and that there is a congregation in the sense of the Principles.

To be sure, the congregation may be present according to its essence, but Sverdrup also wants it to be present according to its form—its proper form.

Hans Nielsen Hauge was a great instrument in the Lord’s hand for leading souls to God. But a congregation according to the demands of the New Testament he could not form, says Sverdrup (C. W. II, 207). The congregation did not receive its proper form in Norway—not through Hauge and not through Johnsen. For Sverdrup reproaches the Johnsen revival for having too little sense for the congregation: “In this whole movement there was lacking the right view of the congregation .... or of Jesus Christ’s body” (C. W. I, 200).

Thus: according to Sverdrup’s view, the congregation’s proper form does not come of itself from what we commonly understand by awakening and conversion.

In all church and Free Church work the salvation of the individual is surely a main matter. But no society would presume to claim that this is its distinctive function. A function that is common to many or to all is distinctive for none.

The principles of the Free Church are distinctive; and the task of the Free Church, which is first mentioned in the Rules (not in the Principles), is distinctive, just as its function is.

The most important thing for us all—and of this we are surely heartily agreed—is the salvation of the individual. But the Free Church believes that the work for the individual’s salvation and preservation in the life of faith, and for spiritual activity, may be better promoted by means of the realization of the Principles, especially Principle No. 1, that principle of which S. said: “The Free Church’s proper main principle is this, that the congregation is the right divine order within the Christian religion” (Guide 143).

I will not say that the editors deny the legitimacy of the local congregation. Quite the contrary. But they do not let the Augsburg understanding of the congregation come clearly forth. The congregation becomes something indefinite, more like the holy catholic church or the communion of saints. And congregational work becomes virtually a mission work, going out to bring souls to life in God.

Now I will not in these observations express my opinion as to whether the line of thought the editors seem to represent is less biblical or more biblical than that represented by Principle No. 1. I do not here set myself up as judge between right and wrong. But I do think that there is some difference in the line of thought.

My position is this:

First, we must become clear about what the Augsburg understanding of the congregation is. And later there may be occasion to test it in the light of God’s Word. And in this testing by the Word, it is not persons that matter, but the Word alone.

For the present, the following may deserve our consideration:
Is there a congregation where two or three are gathered in Jesus’ name?

Yes, says Luther, for whom the church was an organism of all believers on earth, but no organization.

Yes, says Luther, who held that if all believers on earth died, and he were left alone, then the church would still be there in that one believer.

Yes, says Luther, who never spoke of the enlivening of the congregation and never spoke of dead members of the congregation. For his church was a a religious, spiritual reality, independent of the boundaries of time and space, independent of outward organization and form. The holy catholic church, he said, consisted only of believers: it is invisible to others; they cannot grasp that on earth there is an assembly of believers who are bound together not by rule, right, or law, but by the bonds of freedom, Christ’s Spirit. This assembly, or congregation, is visible only to believers. They perceive this unity, this unity of the Spirit; but they cannot point to anything concrete and say: Here is the church. For the church is the work of the Holy Spirit. Everything outward, whether it is called congregation or church, is not church in the biblical sense. Organizations are indeed psychologically necessary, but they are of a worldly (here not sinful) nature. The church, or God’s congregation, is not bound to time, place, organization. This, at any rate, is what the four-hundredth anniversary of the Reformation has sought to establish concerning Luther. I cannot develop this more fully now.

But no, says Sverdrup: the word about “the two or three” is not really spoken concerning the congregation. “If one would apply the word about two or three to the congregation, then the result would be this: if in a place there were no more than two or three believers, then these were nevertheless many enough to be a congregation” (C. W. II, 196).

No, says Sverdrup: Luther’s word in the Third Article concerning the “whole Christian congregation on earth” must also be true of each individual congregation, if this is really to answer to its name. “The congregation consists of the believers in each place, who there in that place have formed an organization in order to labor in God’s kingdom” (Guide 143).

No, says Sverdrup: a believing missionary in a heathen land is no congregation. And he and an assembly of converted people are no congregation, but only a mission.

And no, says Oftedal: “For Paul and the other New-Testament writers the congregation too does not stand as something loose and indefinite, something invisible and universal—the believers, wherever they might be found, without regard to their outward visible union and assembly; but it stands as a closed circle consisting of men and women separated from the rest of the world by their ordered life together and by the word and the Spirit by which they had been separated” (Guide 74). In another place we read that it belongs to the essence of the church to exist in the form of congregations, not as individual Christian persons (1875).

